\documentclass[journal]{IEEEtran}
% Uses local IEEEtran.cls patched to Computer Modern (no Times pfb needed)

\usepackage{cite}
\usepackage{amsmath,amssymb,amsfonts}
\usepackage{graphicx}
\usepackage{textcomp}
\usepackage{xcolor}
\usepackage{booktabs}
\usepackage{multirow}
\usepackage{hyperref}
\usepackage{cleveref}

% Placeholder command for TODO notes during drafting
\newcommand{\todo}[1]{\textcolor{red}{\textbf{[TODO: #1]}}}

\begin{document}

\title{Systematic Review of Blockchain Utilization Strategies in Institutional Consortium Networks}

\author{
  \IEEEauthorblockN{Tony Eneh}
  \IEEEauthorblockA{
    Networked Systems Laboratory (NSL)\\
    Kumoh National Institute of Technology (KIT)\\
    Gumi, South Korea\\
    \todo{email}
  }
}

\maketitle

\begin{abstract}
Institutional actors---banks, regulators, insurers, public agencies, and enterprise consortia---increasingly adopt permissioned and hybrid blockchain architectures to coordinate shared workflows across trust boundaries.
Despite rapid platform proliferation, the technical design space remains fragmented: strategy choices spanning consensus mechanisms, governance models, privacy layers, and interoperability patterns are often made without comparable empirical evidence.
This paper reports the current status of a PRISMA~2020 systematic review targeting blockchain utilization strategies in institutional consortium networks.
The repository currently contains the review protocol, database-specific search workflow, deduplication pipeline, and preliminary title/abstract screening outputs.
Searches were executed across IEEE~Xplore, Web of Science, arXiv, and OpenAlex (used as a practical replacement for ACM Digital Library and Scopus due access constraints).
These searches identified 1,707 records, of which 1,638 remained after deduplication.
Preliminary title/abstract screening marked 571 records for inclusion, excluded 775, and flagged 292 as uncertain.
A two-stage agent workflow then resolved every uncertain record: a rule-based second-pass triage and an automated full-text review of all openly accessible papers (15 downloaded PDFs and 38 open-access landing pages) yielded 143 promotions to include and 149 exclusions (paywalled-only, low-confidence, or out-of-scope), leaving 714 candidate studies (571 first-pass includes plus 143 promoted from uncertain) carried forward to author-level full-text screening.
Full-text screening of those 714 candidates, data extraction, quality assessment, and evidence synthesis across architecture, governance, consensus, privacy, and interoperability dimensions remain ongoing.
\end{abstract}

\begin{IEEEkeywords}
blockchain, consortium networks, permissioned ledger, systematic review, PRISMA, institutional workflows, interoperability, governance, consensus, privacy
\end{IEEEkeywords}

%% ============================================================
%%  INTRODUCTION
%% ============================================================
\section{Introduction}
\label{sec:introduction}

\IEEEPARstart{D}{istributed} ledger technology has evolved well beyond its origins in public cryptocurrency networks.
A growing body of engineering work targets \emph{institutional consortium networks}---multi-party deployments in which banks, regulators, insurers, government agencies, and supply-chain operators share a common ledger infrastructure to coordinate workflows that span organizational trust boundaries~\cite{xu2017taxonomy,belotti2019vademecum,casino2019systematic}.
Unlike public blockchains that optimize for open participation and censorship resistance, consortium architectures prioritize controlled membership, regulatory compliance, deterministic finality, and confidential transaction processing~\cite{hyperledger2018arch,cachin2016architecture}.
Major platform families---Hyperledger Fabric, R3~Corda, ConsenSys Quorum (formerly JP~Morgan Quorum), and Hyperledger Besu---each embed distinct assumptions about node roles, endorsement policies, data visibility, and smart-contract execution models~\cite{androulaki2018hyperledger,brown2016corda,baliga2018performance}.

The promise of consortium blockchains for institutional settings is substantial.
Shared immutable records can reduce reconciliation costs in inter-bank settlement~\cite{guo2016blockchain}, trade finance~\cite{chang2020supplychain}, and cross-border payment corridors~\cite{danezis2015centrally}.
Programmable governance encoded in smart contracts can automate regulatory reporting, enforce access-control policies, and make audit trails tamper-evident~\cite{defilippi2018blockchain}.
In the public sector, multi-agency ledgers have been explored for land registries, digital identity, and health-data exchange where no single entity is trusted to maintain the canonical record~\cite{ølnes2017blockchain,mcghin2019blockchain}.

Despite these advances, the design space is fragmented along several critical dimensions, and practitioners face strategy choices that are poorly supported by comparable evidence.

\subsection{Fragmentation of the Design Space}

First, \emph{consensus and finality strategies} vary widely.
Crash-fault-tolerant protocols such as Raft offer high throughput but assume benign failures, while Byzantine-fault-tolerant protocols such as PBFT, Istanbul~BFT, and newer optimistic or DAG-based variants provide stronger resilience guarantees at the cost of communication complexity and latency~\cite{castro1999practical,sousa2018byzantine,baudet2019state}.
Choosing the right consensus mechanism for a given consortium size, failure model, and throughput target remains an open engineering problem, because benchmarks are conducted under heterogeneous workloads, network topologies, and hardware configurations that resist direct comparison~\cite{dinh2017blockbench,nasir2022performance}.

Second, \emph{governance models} define who can join the network, how policies are updated, and how disputes are resolved.
These models range from static configuration files managed by a founding operator, through multi-signature policy channels, to on-chain decentralized-autonomous-organization (DAO) patterns adapted for permissioned contexts~\cite{beck2018governance,reijers2021governance}.
The interplay between governance agility and security---for example, the risk surface introduced by runtime policy updates versus the operational rigidity of static configurations---is rarely evaluated empirically.

Third, \emph{privacy and confidentiality mechanisms} introduce additional trade-offs.
Channel-based isolation (Hyperledger Fabric), private transaction managers (Tessera in Quorum), zero-knowledge proof integration, and trusted execution environments (TEEs) each offer different privacy guarantees with different performance penalties and trust assumptions~\cite{androulaki2018hyperledger,benhamouda2020can}.
Institutional deployments in regulated sectors (e.g., banking, healthcare) must satisfy stringent data-residency and confidentiality requirements, yet comparative evaluations of privacy mechanisms under realistic multi-party workloads remain scarce.

Fourth, \emph{interoperability} across heterogeneous ledgers and legacy systems is increasingly critical as institutions participate in multiple consortia simultaneously or must bridge on-chain state with off-chain enterprise systems.
Patterns include cross-chain atomic swaps, relay chains, hash-time-locked contracts (HTLCs), API-mediated gateways, oracle networks, and event-driven messaging buses~\cite{belchior2021survey,schulte2019towards}.
Each pattern imposes different trust assumptions, latency overhead, and failure semantics, yet systematic comparison under institutional constraints---where regulatory latency bounds, audit requirements, and counterparty-risk considerations apply---is limited.

\subsection{The Evidence Gap}

Several surveys and mapping studies have catalogued blockchain platforms, consensus algorithms, or application domains at a high level~\cite{casino2019systematic,wang2019blockchain,ali2020comparative}.
However, these existing reviews exhibit three recurring limitations that motivate the present work.

\emph{First}, many reviews focus on platform features or architectural descriptions rather than on empirical, quantitative evaluations of deployed or prototyped systems.
Strategy recommendations are therefore based on specification sheets rather than measured trade-offs under realistic conditions.

\emph{Second}, the institutional consortium context is often conflated with public-chain or single-organization deployments.
The unique constraints of multi-party trust boundaries, regulatory compliance, and organizational governance are not isolated as first-class variables in synthesis.

\emph{Third}, reproducibility and artifact availability are rarely assessed.
Without open benchmarking configurations, workload definitions, and deployment scripts, the community cannot validate or extend reported findings---a critical shortcoming for a field that aspires to engineering rigor.

\subsection{Research Questions}

To address these gaps, this systematic review is structured around four research questions, framed using the PICOC model (Population, Intervention, Comparison, Outcomes, Context) to ensure scope precision:

\begin{enumerate}
  \item[\textbf{RQ1.}] Which blockchain utilization strategies are reported for institutional consortium workflows, and how can they be taxonomized across architecture, governance, consensus, and privacy dimensions?
  \item[\textbf{RQ2.}] What empirical evidence exists on trade-offs among performance, resilience, privacy, and operational complexity for these strategies?
  \item[\textbf{RQ3.}] Which interoperability patterns (cross-chain, API-mediated, bridge/oracle, messaging/event-driven) are most effective in institutional settings, and under what constraints?
  \item[\textbf{RQ4.}] What design gaps and reproducibility weaknesses remain, and how should a reference implementation be structured to address them?
\end{enumerate}

The population comprises institutional multi-party workflows and consortium networks, including inter-bank, inter-agency, and regulated business-to-business collaborations.
The interventions are blockchain utilization strategies encompassing platform choice, consensus protocol, governance model, privacy layer, and interoperability approach.
Comparisons include alternative strategy choices and, where reported, non-blockchain baselines.
Outcomes span throughput, latency, finality, fault tolerance, privacy and security guarantees, interoperability performance, operational overhead, and maintainability.
The context is restricted to regulated or compliance-sensitive settings with multi-organization trust boundaries.

\subsection{Contributions}

This review makes the following contributions:

\begin{itemize}
  \item A \emph{multi-dimensional taxonomy} of blockchain utilization strategies observed in institutional consortium deployments, organized across architecture, governance, consensus, privacy, and interoperability dimensions.
  \item A \emph{structured synthesis of empirical evidence} on performance, resilience, and operational trade-offs, enabling practitioners to compare strategies under stated conditions.
  \item An \emph{assessment of interoperability patterns} specific to institutional constraints, identifying which patterns have been validated and which remain under-explored.
  \item A \emph{reproducibility audit and gap map} that catalogs open weaknesses in the literature and derives requirements for a reference implementation to anchor future experimental work.
\end{itemize}

\subsection{Paper Organization}

The remainder of this paper is organized as follows.
\Cref{sec:background} provides background on consortium blockchain architectures and positions this work relative to prior surveys.
\Cref{sec:method} details the systematic review methodology following PRISMA~2020 and PRISMA-S guidelines.
\Cref{sec:results} presents the results, including study selection, taxonomy construction, and evidence synthesis.
\Cref{sec:discussion} discusses findings, implications for practitioners, and threats to validity.
\Cref{sec:conclusion} concludes the paper and outlines future work toward a reference implementation.

%% ============================================================
%%  BACKGROUND AND RELATED WORK
%% ============================================================
\section{Background and Related Work}
\label{sec:background}

\todo{Subsections: (1) Consortium blockchain architectures (Fabric, Corda, Besu, Quorum); (2) Consensus in permissioned settings; (3) Privacy mechanisms; (4) Interoperability landscape; (5) Positioning vs.\ existing surveys and gap justification.}

%% ============================================================
%%  METHODOLOGY
%% ============================================================
\section{Methodology}
\label{sec:method}

This review follows the PRISMA~2020 reporting guidelines~\cite{page2021prisma} and PRISMA-S extension for search reporting~\cite{rethlefsen2021prismas}.
The protocol draft has been prepared, but Open Science Framework (OSF) registration is still pending.

\subsection{Eligibility Criteria}

Studies were included if they (1)~described a blockchain deployment or prototype in an institutional multi-party setting (inter-bank, inter-agency, or consortium B2B), (2)~provided explicit strategy details covering at least one of consensus, governance, privacy, or interoperability, and (3)~reported quantitative or reproducible technical evaluation outcomes.

Studies were excluded if they were non-technical commentary or policy-only analyses, lacked implementation details, addressed only single-organization private workflows without inter-institutional interaction, or focused exclusively on cryptocurrency market prediction.

\subsection{Information Sources and Search Strategy}

The current repository snapshot searched four operational sources: IEEE~Xplore, Web of Science Core Collection, arXiv, and OpenAlex.
OpenAlex was used as a practical replacement source for ACM Digital Library and Scopus because ACM export access and Scopus search access were unavailable in the working environment.
The search query was constructed from four Boolean blocks addressing technology terms, institutional setting, strategy dimensions, and technical evidence indicators.
Search execution artifacts were frozen on 2026-04-14.
The resulting source counts were 386 records from IEEE~Xplore, 50 from Web of Science, 49 from arXiv, and 1,222 from OpenAlex, for 1,707 records before deduplication.

\subsection{Selection Process}

Records were exported as CSV, deduplicated by DOI, URL, and normalized title-year matching, and managed using a custom deduplication pipeline.
The current repository state includes completion of deduplication and a first-pass title/abstract screening stage, followed by planned full-text assessment against the eligibility criteria.
The first-pass screening was generated by a rule-based triage script aligned to the protocol criteria, assigning each record to INCLUDE, EXCLUDE, or UNCERTAIN classes.
The repository now also includes generated full-text screening and data-extraction sheets for the 863 candidate studies carried forward from the first-pass screen.
The 292 uncertain records were resolved by an agent-driven, auditable two-stage workflow: (i)~a rule-based second-pass triage classified each record as likely\_include, likely\_exclude, or needs\_manual; (ii)~for the needs\_manual subset an open-access discovery script queried OpenAlex and arXiv, downloaded all retrievable open PDFs, and an agent full-text reviewer applied the inclusion criteria to every accessible paper.
The finalisation step then promoted second-pass likely\_include records and agent full-text includes to the candidate pool, and conservatively excluded records that were either (a)~judged likely\_exclude in the second pass, (b)~assessed by the agent as failing the protocol on the available text, (c)~returned needs\_human\_check by the agent (defaulted to exclude), or (d)~had no open full-text access.
Final author-level full-text screening, conflict adjudication, and reviewer statistics for the resulting candidate pool remain ongoing.

\subsection{Data Extraction}

Data items extracted from each included study comprised: bibliographic metadata; platform and technology stack; network model and node roles; consensus and finality strategy; governance and access-control model; privacy and confidentiality mechanism; interoperability strategy; evaluation setup, workload characteristics, and reported metrics; and reproducibility artifacts (code, data, configuration availability).

\subsection{Quality Assessment}

Study quality was assessed across four dimensions: construct validity (alignment of metrics with claims), internal validity (control of confounders), external validity (generalizability to institutional contexts), and reproducibility (artifact completeness and setup determinism).

\todo{Insert scoring rubric details and thresholds.}

\subsection{Synthesis Plan}

Evidence was synthesized through (1)~taxonomy construction across the five strategy dimensions, (2)~quantitative narrative synthesis of metric trade-offs, (3)~identification of evidence concentration and blind spots, and (4)~derivation of requirements for a reference implementation.

%% ============================================================
%%  RESULTS
%% ============================================================
\section{Results}
\label{sec:results}

\subsection{Study Selection}

The current study-selection pipeline has completed the identification, deduplication, and preliminary title/abstract screening phases.
Across the four operational sources, 1,707 records were identified: 386 from IEEE~Xplore, 50 from Web of Science, 49 from arXiv, and 1,222 from OpenAlex.
Deduplication removed 69 records (40 DOI matches and 29 title-year matches), leaving 1,638 unique records for screening.
Preliminary title/abstract screening excluded 775 records, marked 571 as included, and flagged 292 as uncertain for adjudication.
The agent-driven uncertain-resolution workflow then processed all 292 uncertain records: 123 were promoted to include via second-pass triage, 20 additional records were promoted via accessible full-text agent review, 1 was excluded after full-text review, 32 were conservatively excluded after needs\_human\_check verdicts, 48 were excluded as second-pass likely\_exclude, and 68 were excluded for lack of open full-text access (paywalled-only).
The finalised uncertain partition is therefore 143 included and 149 excluded, and the effective candidate pool for author-level full-text review is now 714 records (571 first-pass includes plus 143 promoted from uncertain).

\todo{Insert PRISMA flow diagram figure once full-text review and exclusion-reason coding are complete.}

\subsection{Study Characteristics}

Study-characteristics extraction has not yet been completed because the final full-text inclusion set is still under construction.
An extraction sheet has already been generated in the repository for the 863 candidate studies that survived first-pass screening, covering publication year, venue, domain, platform, consensus, governance, privacy mechanism, interoperability approach, workload, metrics, and reproducibility artifacts; the 714-record post-uncertain-finalisation pool is a strict subset of that sheet.

\subsection{RQ1: Taxonomy of Blockchain Utilization Strategies}

\todo{Present multi-dimensional taxonomy across architecture, governance, consensus, privacy, and interoperability. Include classification table and discussion of dominant patterns.}

\subsection{RQ2: Empirical Trade-Offs}

\todo{Summarize quantitative evidence on throughput, latency, finality, fault tolerance, privacy overhead, and operational complexity. Comparative tables grouped by strategy dimension.}

\subsection{RQ3: Interoperability Patterns}

\todo{Categorize cross-chain, API-mediated, bridge/oracle, and messaging/event-driven patterns. Report effectiveness metrics and institutional constraints.}

\subsection{RQ4: Gaps and Reference Implementation Requirements}

\todo{Reproducibility audit results. Gap map. Derived requirements for reference implementation.}

%% ============================================================
%%  DISCUSSION
%% ============================================================
\section{Discussion}
\label{sec:discussion}

\todo{Subsections: (1) Principal findings and their implications; (2) Practical recommendations for institutional adopters; (3) Comparison with prior surveys; (4) Threats to validity (search completeness, selection bias, synthesis limitations).}

%% ============================================================
%%  CONCLUSION
%% ============================================================
\section{Conclusion}
\label{sec:conclusion}

\todo{Summarize contributions, key findings per RQ, limitations, and future work (reference implementation plan).}

%% ============================================================
%%  REFERENCES
%% ============================================================
\bibliographystyle{IEEEtran}
\bibliography{references}

\end{document}
